% Brief (grammar: linear provenance chain, time left to right; atlas
%   V/fig:anchor-delegation-inline).
% Reader question: what evidence travels with a multi-hop delegated
%   capability, and what happens when an attacker splices a hop?
% Claim: each signed hop binds parent, child and a fresh nonce to one shared
%   task hash h; a spliced-in signature that does not cover the local nonce
%   breaks exactly at the one crossing it touches.
% Evidence: identities id1-id4 (\S "Phase 3: Stigmergic Delegation"); ordinary
%   hops 1 and 3 (sig1, sig3, fresh n1/n3); the splice at hop 2 (sig3 moved
%   into the id2-id3 slot, where a signature over n2 is expected) [internal].
% Must distinguish: an ordinary hop (focus, solid) from the spliced hop
%   (breach, dashed) -- both travel the same adversary-visible network.
% Grammar chosen: tree/DAG left to right (tpl-tree-dag), one dashed network
%   band behind all three crossings per chapter line ~746 (added here; it was
%   missing before this redraw, flagged for the author to confirm). Rejected:
%   a generic node-link graph -- the atlas's own rejection, since it would
%   erase the left-to-right time order the splice depends on.
% Five-second test: "three crossings, same shared hash h below; the middle
%   one is dashed red because its signature does not cover n2."
\begin{figure}[H]
\centering
\pdfigurehue{pdcobalt}%
\begin{tikzpicture}[pd figure,x=.66cm,y=1cm]
  \draw[pd warn rule] (-5.75,1.45) rectangle (5.75,3.35);
  \node[pd label,anchor=south west] at (-5.7,3.9) {adversary-visible};
  \foreach \x/\id in {-5.1/1,-1.7/2,1.7/3,5.1/4} {
    \node[pd focus outline,circle,minimum size=9mm,inner sep=0pt] (id\id) at (\x,2.1) {id\id};
  }

  % ---- crossing 1: id1 -> id2, ordinary hop --------------------------------
  \draw[pd focus arrow] (id1.east) to[out=15,in=165] (id2.west);
  \node[pd focus label,anchor=south] at (-3.4,2.42) {sig1};
  \node[pd focus datum,minimum size=5pt,inner sep=0pt] (nonce1) at (-3.4,.65) {};
  \node[pd label,anchor=west,align=left,text width=1.5cm] at (-3.26,.65) {n1 fresh};
  \draw[pd focus rule] (-3.4,1.65) -- (nonce1);

  % ---- crossing 2: id2 -> id3, the attempted splice ------------------------
  \draw[pd breach rule] (id2.east) to[out=15,in=165] (id3.west);
  \node[pd breach datum] (splice) at (0,2.1) {};
  \node[pd label,text=pderror,align=center,text width=3.5cm,anchor=south]
    at (0,2.85) {sig3 does not cover n2: refused};
  \node[pd breach datum] (nonce2) at (0,.65) {};
  \node[pd label,anchor=west,align=left,text width=1.5cm] at (.14,.65) {n2 expected};
  \draw[pd breach rule] (0,1.65) -- (nonce2);

  % ---- crossing 3: id3 -> id4, ordinary hop --------------------------------
  \draw[pd focus arrow] (id3.east) to[out=15,in=165] (id4.west);
  \node[pd focus label,anchor=south] at (3.4,2.42) {sig3};
  \node[pd focus datum,minimum size=5pt,inner sep=0pt] (nonce3) at (3.4,.65) {};
  \node[pd label,anchor=west,align=left,text width=1.5cm] at (3.54,.65) {n3 fresh};
  \draw[pd focus rule] (3.4,1.65) -- (nonce3);

  \draw[pd rule] (-5.55,-.35) -- (5.55,-.35);
  \foreach \x in {-5.1,-1.7,1.7,5.1} \node[pd focus datum,minimum size=4pt,inner sep=0pt] at (\x,-.35) {};
  \node[pd label,anchor=north west,text width=1.9cm,align=left] at (-5.55,-.55) {task hash h};

  \node[pd note,anchor=north west,align=left,text width=6.9cm] at (-5.55,-1.15)
    {each ordinary crossing binds parent, child, a fresh nonce, and the same
     task hash h; the breach crossing reuses sig3 (bound to n3) where n2 was
     expected, so verification fails at that crossing only};
\end{tikzpicture}
\caption{The delegation chain as a signed provenance braid across four
identities id1--id4, drawn inside the network the Dolev-Yao adversary
observes. Ordinary hops 1 and 3 bind parent, child, a fresh nonce (n1, n3),
and the shared task hash h below. Hop 2 shows a splice attack: an attacker
moves the legitimate signature sig3 (computed over id3, id4, and n3) into
the id2--id3 position, where the verifier expects a signature over n2; the
check fails at exactly that crossing, so the discontinuity is visible at one
adjacent-identity boundary rather than anywhere else in the chain.
[internal]}
\label{fig:anchor-delegation-inline}
\end{figure}
